% Example: Polish official letter -- package pl-official-letter
% Build (from repo root, TEXINPUTS points at templates/):
%   TEXINPUTS=./templates//: latexmk -pdf -interaction=nonstopmode -halt-on-error -outdir=out examples/official-letter-example.tex
\documentclass[11pt,a4paper]{article}
\usepackage{pl-official-letter}
% Optional: how much lower the sender starts vs. place line (default 2.6em)
% \plnadawcatopgap{3.2em}
% Space between address block and subject title (default 2.75em)
% \plheadtotitlespace{3.25em}
% Address block column widths (defaults: left 0.46, right 0.48 \textwidth)
% \plleftcolwidth{0.47\textwidth}
% \plrightcolwidth{0.47\textwidth}
% Gap from date line to recipient (right column)
% \plsendertorecipient{1.4em}

% PDF metadata (optional)
\plpdftitle{Przykładowy wniosek do urzędu}
\plpdfauthor{Jan Kowalski}

\begin{document}

% --- Sender (multi-line: \\) ---
\plnadawca{%
  Jan Kowalski\\
  ul. Przykładowa 1 m. 2\\
  00-001 Warszawa%
}

% --- Recipient: in PDF on the right, under place and date ---
\pladresat{%
  Urząd Gminy Przykład\\
  Wydział Spraw Obywatelskich\\
  ul. Urzędowa 10\\
  00-999 Przykładowo%
}

% --- Place + date (omit word "dnia" and "r." -- the template adds them) ---
\plmiejscodata{Warszawa}{16 kwietnia 2026}
% Alternatively: \plmiejscowosc{...} and \pldata{...}

% --- Subject title (centered, bold) ---
\pltytul{Wniosek o zmianę adresu zameldowania na pobyt stały}

% --- Optional: case reference / scope (remove entire \plsprawa{} line if unused) ---
\plsprawa{znak sprawy URG-WSC/123/2026}

% --- Salutation (package default unless overridden with \plzwrot{...}) ---
% \plzwrot{Szanowny Panie Burmistrzu,}
% Optional hint marking the handwritten signature area:
% \plpodpowiedzpodpisu{(handwritten signature)}

\makeletterhead

W związku z nabyciem prawa do lokalu mieszkalnego przy ul. Nowej 5 w Przykładowie zwracam się z prośbą o dokonanie wpisu zmiany adresu zameldowania na pobyt stały.

Do wniosku dołączam wymagane prawem załączniki. Proszę o rozpatrzenie sprawy w terminie określonym przepisami.

\plpodpis{Jan Kowalski}
% Optional line under the name:
% \plfunkcja{handwritten signature}

\makelettersignature

\end{document}
